%% S1.5: Productive terms of S.  Neutral amsart shell, as S1; port to the venue class later.
%%
%% Paper-to-Lean map (lean-ski v1.5.0; all in the #audit_clean gate, Ski/Audit.lean).
%%   head reduction ->h*, hnf of t ...... HnfSound.HeadChain, sound_head, headChain_uniq
%%   P_k, productive ................... HnfProd.Pk; HnfCert.Prod
%%   gamma, c(t), class 37 ............. HnfProd.clsTable, HnfProd.γ, HnfProd.cls
%%   lem:table (i)-(iv) ................ HnfProd.slaw (slawB_eq), absorbL/absorbR (absorbB_eq),
%%                                       pair_mem (startB_eq, pairsClosedB_eq, reachBitsB_eq), b1 (b1B_eq)
%%   lem:invc ........................... HnfProd.wsteps_cls
%%   prop:easy .......................... HnfProd.easy
%%   lem:b0 ............................. HnfB0.B0
%%   lem:ci ............................. HnfQ5.pk_wsteps (items_join, transport)
%%   lem:sub ............................ HnfQ5.prod_sub
%%   (A), prop:reduce ................... HnfQ5.MinimalPairs, hard_of_A, q5_of_A
%%   lem:laws F', L0, L1, LA ............ HnfCertLaws.chain_factor, L0, L1, LA
%%   exact sizes, minimum sizes ......... HnfCert.exact_size (exactOK_all), minsz_le (MinOK)
%%   shape table, split table ........... HnfCert.unf_s0/s1/s2 (ShapesOK), splits_complete (SplitsOK)
%%   instances .......................... HnfCert.Inst, InstR
%%   the checker ........................ HnfCert.checkClaim, run, hnfStep, pairClose, grade
%%   lem:cert (soundness) ............... HnfCert.claims_sound, run_sound, split_sound
%%   coverage ........................... HnfCert.rootsCovered (rootsCovered_eq)
%%   prop:A ............................. HnfCert.minimalPairs (HnfCertFinal.lean)
%%   thm:main ........................... HnfCert.q5
%%   certificate data ................... Ski/HnfCertData.lean (generated by
%%                                       notes/checks/s15/cert/klean_gen.py in the ski tree)
\documentclass[11pt]{amsart}
\usepackage[T1]{fontenc}
\usepackage{lmodern}
\usepackage{amsmath,amssymb,amsthm}
\usepackage[numbers,sort]{natbib}
\usepackage{booktabs}
\usepackage[hidelinks]{hyperref}

\newtheorem{theorem}{Theorem}
\newtheorem{lemma}[theorem]{Lemma}
\newtheorem{proposition}[theorem]{Proposition}
\newtheorem{corollary}[theorem]{Corollary}
\theoremstyle{definition}
\newtheorem{definition}[theorem]{Definition}
\theoremstyle{remark}
\newtheorem{remark}[theorem]{Remark}

\newcommand{\spine}{\operatorname{fr}}
\newcommand{\lean}[1]{\texttt{#1}}
\newcommand{\app}{\mathbin{+\!\!+}}
\newcommand{\hred}{\to_{h}^{*}}
\newcommand{\rej}{\mathsf{r}}
\newcommand{\sz}[1]{|#1|}
\newcommand{\st}{\operatorname{st}}

\title{Productive Terms of the Combinator $S$}
\author{N.\ E.\ Davis}
\address{ZeroAttest, Champaign, Illinois, USA}
\email{NDavis@ZeroAttest.com}
\urladdr{https://orcid.org/0000-0002-1131-8319}
\date{\today}

\begin{document}
\begin{abstract}
A closed term built from $S$ alone is productive when its B\"ohm tree has no $\bot$.
Barendregt, Endrullis, Klop and Waldmann asked how to characterize these terms, whether
they form a regular set, and whether productivity is decidable (\emph{Dance of the
Starlings}, Question~5). We give a 39-state bottom-up tree automaton that accepts
exactly the productive closed $S$-terms, so the set is regular and productivity is
decidable. The automaton is a model of the $S$-law, so its state is a reduction
invariant, and acceptance implies productivity by the head-normal-form theorem of a
companion paper. The converse reduces to 2,757 pairs of states, which a certificate of
7,244 claims about symbolic head reduction settles. The proof, including the soundness of
the certificate checker and its run on the certificate, is mechanized in Lean~4 with no
axioms beyond propositional extensionality and quotient soundness, for productivity stated
as ``the B\"ohm tree is defined to every depth''; its equivalence with the B\"ohm-tree
statement is a short argument on paper.
\end{abstract}

\keywords{combinatory logic, combinator $S$, B\"ohm trees, productivity, tree automata,
certificates, Lean}
\maketitle

\section{Introduction}\label{sec:intro}

\citet{dots} call a term \emph{productive} when its B\"ohm tree has no $\bot$. For example $S\,T\,T\,(S\,T\,T)$ with $T=S\,S$ is productive
with an infinite B\"ohm tree, while $A\,A\,A$ with $A=S\,S\,S$ is not. Their
Question~5 asks how to characterize the productive closed $S$-terms, whether they form a
regular set, and whether productivity is decidable. For normalization the answers are
known: Waldmann showed that the normalizing closed $S$-terms form a regular tree language
and that normalization is decidable~\citep{waldmann-thesis,waldmann-ic}; see
\citet[Theorem~9]{dots}. Productivity is weaker than normalization, since a productive term
may have an infinite B\"ohm tree. In full combinatory logic productivity is undecidable, by
the Scott--Curry theorem, since it is non-trivial and closed under
conversion~\citep{barendregt}. The companion paper~\citep{s1}
settles the first layer of the B\"ohm tree: a 38-element model of the $S$-law decides
whether a term has a head normal form, confirming the automaton of their Question~10. Here we
settle the whole tree.

\begin{theorem}\label{thm:main}
There is a bottom-up tree automaton with 39 states, one of them rejecting, that accepts a
closed $S$-term if and only if the term is productive. Hence the productive closed
$S$-terms form a regular tree language, and productivity is decidable.
\end{theorem}

The automaton's transition table $\gamma$ is a model of the $S$-law (Section~\ref{sec:model}),
so the state of a term is invariant under reduction, and an accepted term is
productive by the soundness theorem of~\citep{s1}. The converse needs more. A rejected term
contains a least rejected subterm $u\,v$, with $u$ and $v$ accepted, and non-productivity
passes from subterms to terms (Section~\ref{sec:reduce}); so it suffices that every such
\emph{minimal pair} is non-productive. That is a statement about infinitely many terms
but only 2,757 pairs of states, and Section~\ref{sec:cert} proves it with a certificate:
claims of the form ``every instance of this symbolic term is non-productive'', each
checked by a short symbolic head reduction or a case split and justified by smaller claims. Instance sizes
make the justification well founded, and three laws about how head reduction changes
sizes (Section~\ref{sec:laws}) supply the size bounds.

The proof, including the run of the checker on the certificate, is mechanized in
Lean~4~\citep{lean4} within \texttt{lean-ski}~\citep{leanski}. We use the notation of~\citep{s1} and~\citep{barendregt}.

\section{Terms, head reduction and B\"ohm trees}\label{sec:prelims}

We write $\mathcal S$ for the closed terms built from $S$ by application, $\to$ for weak
reduction and $\twoheadrightarrow$ for its reflexive--transitive closure. Every
$t\in\mathcal S$ is uniquely $S\,u_1\cdots u_m$; the list $\spine(t)=[u_1,\dots,u_m]$ is
its \emph{spine}. A term is a \emph{head normal form} (hnf) if its spine has at most two
items. If $\spine(t)=[x,y,z]\app\rho$, the \emph{head step} is
$t\to_h x\,z\,(y\,z)\,\rho$, with
\begin{equation}\label{eq:headstep}
\spine(x\,z\,(y\,z)\,\rho)=\spine(x)\app[z,\;y\,z]\app\rho .
\end{equation}
Head reduction $\hred$ is deterministic and stops at an hnf; so if $t$
head-reduces to an hnf, that hnf is unique, and we call it $h(t)$. The \emph{size}
$\sz t$ is the number of occurrences of $S$ in~$t$.

We use from~\citep{s1} the block model: a table $\beta$ on $B=\{0,\dots,37\}$ satisfying the
$S$-law, the fold $b(S)=0$, $b(u\,v)=\beta(b(u),b(v))$, invariant under reduction, and the
left-absorbing block $\bot=36$: $\beta(\bot,a)=\bot$ for all $a$.

\begin{theorem}[\citep{s1}]\label{thm:s1}
For $t\in\mathcal S$, $b(t)\ne\bot$ if and only if $t$ head-reduces to an hnf; and if
$b(t)=\bot$ then $t$ has no hnf at all.
\end{theorem}

\begin{definition}
Let $P_0=\mathcal S$, and let $t\in P_{k+1}$ when $t$ head-reduces to an hnf and every item
of $\spine(h(t))$ is in $P_k$. A term is \emph{productive} if it lies in every~$P_k$.
\end{definition}

The \emph{B\"ohm tree} $\mathrm{BT}(t)$ is $\bot$ if $t$ has no hnf, and otherwise
$S\,\mathrm{BT}(a_1)\cdots\mathrm{BT}(a_j)$ for an hnf $S\,a_1\cdots a_j$ of $t$; it does not
depend on the hnf chosen, since two hnfs of $t$ have a common reduct, so they have the same
number of items and corresponding items are convertible~\citep{barendregt}. By
Theorem~\ref{thm:s1}, $t$ has an hnf iff it head-reduces to one, so we may take $h(t)$. Then,
by induction on $k$, $t\in P_k$ iff $\mathrm{BT}(t)$ has no $\bot$ at depth less than $k$:
$t\in P_{k+1}$ iff the root of $\mathrm{BT}(t)$ is not $\bot$ and the items of $h(t)$, whose
trees are the subtrees at depth $1$, lie in $P_k$. Hence $t$ is productive in this sense iff
$\mathrm{BT}(t)$ has no $\bot$, the definition of~\citep{dots}. Lemma~\ref{lem:ci} shows that
$P_k$ is invariant under reduction.

\section{The class model and the easy half}\label{sec:model}

Let $C=\{0,\dots,38\}$ and let $\gamma\colon C\times C\to C$ be the table of
\lean{HnfProd.clsTable}, also supplied as the file \texttt{gamma.csv} with this paper. It is
not a refinement of $\beta$: the class of a term does not determine its block, and the pair
takes 115 values (Lemma~\ref{lem:table}(iii)). The elements of $C$ are \emph{classes}; class $0$ interprets $S$, and we
write $\rej=37$. The \emph{class} of a term is the fold $c(S)=0$,
$c(u\,v)=\gamma(c(u),c(v))$, and a term is \emph{accepted} if $c(t)\ne\rej$. The automaton
of Theorem~\ref{thm:main} has states $C$, leaf state $0$, transitions $\gamma$ and
accepting states $C\setminus\{\rej\}$. The \emph{product state} of $t$ is
$\st(t)=(c(t),b(t))$.

\begin{lemma}[table facts]\label{lem:table}
For all $a,a',a''\in C$:
\begin{enumerate}
\item[(i)] $\gamma(\gamma(\gamma(0,a),a'),a'')=\gamma(\gamma(a,a''),\gamma(a',a''))$
  \textup{($39^3=59{,}319$ instances)};
\item[(ii)] $\gamma(\rej,a)=\gamma(a,\rej)=\rej$;
\item[(iii)] the set $R$ of 115 product states in \lean{HnfProd.reachPairs} contains $(0,0)$
  and is closed under $(\gamma,\beta)$, so $\st(t)\in R$ for every $t\in\mathcal S$;
\item[(iv)] for $(c,b)\in R$, $c\ne\rej$ implies $b\ne\bot$.
\end{enumerate}
\end{lemma}

\begin{lemma}[invariance]\label{lem:invc}
If $t\twoheadrightarrow t'$ then $c(t')=c(t)$.
\end{lemma}
\begin{proof}
As for blocks in~\citep{s1}: a contraction is Lemma~\ref{lem:table}(i) at the classes of
$x,y,z$, and the fold is compositional.
\end{proof}

\begin{proposition}[easy half]\label{prop:easy}
Every accepted $t\in\mathcal S$ is productive.
\end{proposition}
\begin{proof}
We show $t\in P_k$ by induction on $k$, for all accepted $t$ at once. By
Lemma~\ref{lem:table}(iii,iv), $b(t)\ne\bot$, so $t$ head-reduces to an hnf $h(t)$ by
Theorem~\ref{thm:s1}, and $c(h(t))=c(t)\ne\rej$ by Lemma~\ref{lem:invc}. Since $\rej$ is
absorbing on both sides (Lemma~\ref{lem:table}(ii)), the items of $h(t)$ are accepted, and
by induction they lie in~$P_k$.
\end{proof}

\section{Reduction to minimal pairs}\label{sec:reduce}

\begin{lemma}[B0]\label{lem:b0}
If some subterm of $t\in\mathcal S$ has block $\bot$, then $t$ is not productive.
\end{lemma}
\begin{proof}
Say that $t$ \emph{has a hole within $n$} if $b(t)=\bot$, or $n>0$ and some item of
$\spine(t)$ has a hole within $n-1$. A subterm at any position is reached through items of
spines, and a hole in the function part of an application stays a hole of the application
(by absorption, $\beta(\bot,a)=\bot$), so $t$ has a hole within some~$n$. A head step
preserves ``hole within $n$'': if $b(t)=\bot$, so is the reduct's block, since the block is
invariant~\citep{s1}; by~\eqref{eq:headstep}, $z$ and the items of $\rho$ stay items; a hole
in $y$ is a hole in $y\,z$; a hole in an item of $x$ becomes a hole in an item of the reduct
(a hole within $n-2$ is one within $n-1$); and if $b(x)=\bot$ then the reduct has block $\bot$. Now suppose $t\in P_{n+1}$.
Then $h(t)$ has a hole within~$n$, and its block is not $\bot$ (Theorem~\ref{thm:s1}), so
for $n=0$ this is impossible, and for $n>0$ some item of $h(t)$ has a hole within $n-1$ and
lies in $P_n$. By induction on $n$ that is impossible too, so $t\notin P_{n+1}$.
\end{proof}

\begin{lemma}[reduction invariance]\label{lem:ci}
If $t\in\mathcal S$ and $t\twoheadrightarrow t'$, then $t\in P_k\iff t'\in P_k$, for every~$k$.
\end{lemma}
\begin{proof}
By induction on $k$, for all such pairs. Suppose that one of $t,t'$ lies in $P_{k+1}$. That term
has an hnf, so $b(t)=b(t')\ne\bot$ and both head-reduce to hnfs $h(t)$ and $h(t')$
(Theorem~\ref{thm:s1}). These are convertible, so by Church--Rosser they have a common
reduct. A reduct of $S$, $S\,a$ or $S\,a\,b$ is $S$, $S\,a^*$ or $S\,a^*\,b^*$ with
$a\twoheadrightarrow a^*$ and $b\twoheadrightarrow b^*$, so the two hnfs have the same
number of items and corresponding items have a common reduct. By the induction hypothesis
corresponding items agree on $P_k$, so both terms lie in $P_{k+1}$.
\end{proof}

\begin{lemma}[subterms]\label{lem:sub}
Every subterm of a productive $t\in\mathcal S$ is productive.
\end{lemma}
\begin{proof}
We show $m\in P_k$ for every subterm $m$ of every productive $t$, by induction on $k$. By
Lemma~\ref{lem:b0}, $b(m)\ne\bot$, so $m$ head-reduces to $h(m)$. Replacing $m$ by $h(m)$ in
$t$ gives a reduct $t'$ of $t$, productive by Lemma~\ref{lem:ci}, and the items of $h(m)$
are subterms of~$t'$; by induction they lie in~$P_k$, so $m\in P_{k+1}$.
\end{proof}

\begin{proposition}\label{prop:reduce}
Suppose that
\begin{itemize}
\item[\textup{(A)}] for all accepted $u,v\in\mathcal S$ with $c(u\,v)=\rej$ and
  $b(u\,v)\ne\bot$, the term $u\,v$ is not productive.
\end{itemize}
Then every rejected $t\in\mathcal S$ is not productive.
\end{proposition}
\begin{proof}
By induction on $t$. The term $S$ is accepted. Let $t=u\,v$ be rejected. If $u$ or
$v$ is rejected, it is not productive by induction, so $t$ is not productive by
Lemma~\ref{lem:sub}. If $b(t)=\bot$, Lemma~\ref{lem:b0} applies. Otherwise (A) applies.
\end{proof}

Whether (A) applies to $u\,v$ depends only on $\st(u)$ and $\st(v)$, and Lemma~\ref{lem:table}(iii) leaves
95 accepted product states and 2,757 pairs $(\st(u),\st(v))$ to which (A) applies.

\section{Size laws for head reduction}\label{sec:laws}

\begin{lemma}\label{lem:laws}
Let $t\in\mathcal S$.
\begin{enumerate}
\item[(F)] If $t\hred w\,L$, then either $t=t_0\,L$ with $t_0\hred w$, or
  $t\hred S\,x\,y\,z$ with $L=y\,z$ and $x\,z\hred w$.
\item[(L0)] If $t\hred u$ and $\spine(u)$ has at most one item, then $t=u$.
\item[(L1)] If $t\hred S\,S\,y$, then $\sz y=\sz t-2$.
\item[(LA)] If $t\hred S\,(S\,S)\,y$, then $\sz y\le\sz t-2$.
\end{enumerate}
\end{lemma}
\begin{proof}
(F) By induction on the reduction. A head step from a term of arity at least four leaves
the last argument in place; a step from arity exactly three, $S\,x\,y\,z\to x\,z\,(y\,z)$,
makes $y\,z$ the last argument.

(L0) By~\eqref{eq:headstep}, a head step yields a spine of at least two items, and
later steps keep that.

(L1) By induction on $\sz t$. Apply (F) with $L=y$. In the first case $t_0\hred S\,S$, so
$t_0=S\,S$ by (L0) and $\sz y=\sz t-2$. In the second, $x\,z\hred S\,S$ gives $x\,z=S\,S$,
so $x=z=S$, $y=y'\,S$ and $t\hred S\,S\,y'\,S$. The last argument $S$ is an atom, so it
was not created by a step: by (F), $t=t_1\,S$ with $t_1\hred S\,S\,y'$, and by induction
$\sz y=\sz{y'}+1=\sz{t_1}-1=\sz t-2$.

(LA) Apply (F) with $L=y$. In the first case $t_0=S\,(S\,S)$ by (L0), and
$\sz y=\sz t-3$. In the second, $x\,z=S\,(S\,S)$, so $x=S$, $z=S\,S$, $y=y'\,(S\,S)$ and
$t\hred S\,S\,y'\,(S\,S)$. Apply (F) again with $L=S\,S$. If $t=t_1\,(S\,S)$ with
$t_1\hred S\,S\,y'$, then (L1) gives $\sz y=\sz t-2$. Otherwise $S\,S=y_2\,z_2$ was created,
so $y_2=z_2=S$ and $t\hred S\,x_2\,S\,S$ with $x_2\,S\hred S\,S\,y'$; by (L1),
$\sz{y'}=\sz{x_2}-1$. The two atomic arguments were not created, so by (F) twice
$t=t_4\,S\,S$ with $t_4\hred S\,x_2$, hence $t_4=S\,x_2$ by (L0), and
$\sz y=\sz{x_2}+1=\sz t-2$.
\end{proof}

Two finite facts complete the size bounds. First, sixteen blocks, listed in
\lean{HnfCert.EXACTT}, contain terms of one size only (block $1$ is $\{S\,S\}$, for example);
this follows by induction on terms, with a finite check that every $\beta$-product landing in
such a block comes from blocks of fixed sizes that add up. Second, a table $\mu$ of lower bounds on
sizes is subadditive on $R$, $\mu(\st(u\,v))\le\mu(\st(u))+\mu(\st(v))$ with $\mu(\st(S))=1$,
so $\sz t\ge\mu(\st(t))$.

\section{The certificate}\label{sec:cert}

\paragraph{Symbolic terms.} A \emph{symbolic term} is built by application from $S$ and
variables $X_\sigma$, one for each accepted product state $\sigma\ne\st(S)$; we write $X_{\st(S)}$ for $S$. A term $u$ is
an \emph{instance} of $X_\sigma$ if $\st(u)=\sigma$, of $S$ if $u=S$, and of $T_1\,T_2$ if
$u=u_1\,u_2$ with $u_i$ an instance of~$T_i$. A \emph{claim} is a symbolic term $T$ with a
\emph{stratum} and a \emph{rank} in~$\mathbb N$; it is \emph{valid} if no instance of $T$
is productive. Claims are ordered by stratum, then rank; for claims about instances
$u$ and $u'$ we use the lexicographic order on (stratum, size, rank),
$(s',\sz{u'},r')<(s,\sz u,r)$.

\paragraph{Split claims.} A split claim names an occurrence of a variable $X_\sigma$ in
$T$, and is checked by finding, for every pair $(p,q)$ of accepted product states such that
a term of state $p$ applied to a term of state $q$ has state $\sigma$, a claim on $T[X_p\,X_q]$ (the occurrence replaced) that is lower in
(stratum, rank). This is sound: an instance of $X_\sigma$ is an application of two
accepted terms (Lemma~\ref{lem:table}(ii)), so an instance of $T$ is an instance of
some $T[X_p\,X_q]$, of the same size. A table, checked complete, lists the pairs.

\paragraph{Run claims.} A run claim has at most five variable occurrences, of states $\sigma_1,\dots,\sigma_5$, whose
instances have sizes $s_1,\dots,s_5$ (or the fixed size of the occurrence's block). An instance
has size $N$, a linear form in the $s_i$, and $s_i\ge\mu(\sigma_i)$ by the minimum sizes. The check runs head reduction on
$T$ symbolically. Every leaf of the running term carries a size bound, a linear form in the
$s_i$ or ``unknown'', and every node its product state. For a symbolic subterm whose
instances have size at most $f$, the \emph{grade} is $2$ if comparing coefficients shows $f\le N-1$
for all admissible $s_i$, $1$ if it shows $f\le N$, and $0$ otherwise. A claim $d$ \emph{covers} a subterm of grade $g$
if $d$ is in a lower stratum, or in the same stratum with $g=2$, or with $g=1$ and a lower
rank. At each stage the first applicable rule treats the running term $U$.
\begin{enumerate}
\item[(R1)] If a node of $U$ has block $\bot$, the branch closes.
\item[(R2)] If a rejected node has two accepted children with states $p,q$, and a claim on
  $X_p\,X_q$ lies in a lower stratum, or in the same stratum with the node of grade~$2$, the
  branch closes.
\item[(R3)] If $\spine(U)=[x,y,z]\app\rho$ with head $S$, take the head step.
\item[(R4)] If $U$ is an hnf with head $S$, the next \emph{hint}, a number stored with the
  claim, names a rejected item $a$ and one of two moves. Either a claim covers the
  \emph{generalization} of $a$, the symbolic term obtained by replacing each maximal accepted
  proper subterm of $a$ by the variable of its state, at the grade of $a$ (computed from
  the size bound of $a$), and the branch closes; or the run continues on $a$, a
  \emph{descent}.
\item[(R5)] If the head of $U$ is a variable $X_\sigma$ with at least one argument, the run
  continues on every term obtained by replacing $X_\sigma$ with an hnf shape of $\sigma$, namely
  $S\,X_a$ with $\st(S\,X_a)=\sigma$ or $S\,X_a\,X_b$ with
  $\st(S\,X_a\,X_b)=\sigma$, for accepted $a$ and $b$ (a table checked complete lists the
  shapes; the shape $S$ is excluded because $\sigma\ne\st(S)$ and only $S$ has state $\st(S)$). The
  new leaf $X_a$ in $S\,X_a$ has size bound one less than the head's, by (L0). In
  $S\,X_a\,X_b$, the bound of $X_b$ is two less than the head's if $a$ has block $0$ (the
  term $S$) or $1$ (the term $S\,S$), by (L1) and (LA), and unknown otherwise. Leaves in
  blocks of fixed size get that size; the leaf $X_a$ of $S\,X_a\,X_b$ otherwise has no bound.
\end{enumerate}
The run fails if no rule applies or if its fuel (3,000 stages along each branch) or the number of nested descents along a branch exceeds 10.

\begin{lemma}[soundness]\label{lem:cert}
If the tables of minimum sizes, hnf shapes and splits check, and every claim checks, then every
claim is valid.
\end{lemma}
\begin{proof}
By induction on (stratum, size, rank): let $u$ be an instance of a claim $c$, and assume
every claim valid on instances lower in the order. A split claim reduces $u$ to a claim of
lower (stratum, rank) on the same $u$. For a run claim we show, by induction on the remaining
fuel, that if the run succeeds from a running term $U$, then no instance of $U$ that obeys the
bounds and states is productive; $u$ is such an instance of $T$. In (R1) the instance has a
subterm of block $\bot$, and Lemma~\ref{lem:b0} applies. In (R2) the covering claim is lower
in the order at the instance of the node: if it lies in a lower stratum, whatever the size;
if in the same stratum, the node has grade~$2$, so its instance has size at most $N-1$. So
that subterm is not productive, and Lemma~\ref{lem:sub} applies. In (R3) the instance takes the same head step,
and Lemma~\ref{lem:ci} applies. In (R4) the instance of the chosen item is a subterm of the
instance of $U$; it is not productive, either because the covering claim is lower in the order
there (its size is at most $N-1$ for grade~2, or $N$ with a lower rank for grade~1) or by the
induction hypothesis on the continued run; Lemma~\ref{lem:sub} applies. In (R5)
$\st(\text{head})=\sigma$ is accepted, so by Lemma~\ref{lem:table}(iv) and
Theorem~\ref{thm:s1} the head of the instance head-reduces to an hnf of state~$\sigma$, whose
items are accepted by Lemma~\ref{lem:table}(ii); the reduct is an instance of one of the
listed shapes, with the bounds of Lemma~\ref{lem:laws} and the fixed sizes, and
Lemma~\ref{lem:ci} applies.
\end{proof}

\begin{proposition}\label{prop:A}
Hypothesis \textup{(A)} holds.
\end{proposition}
\begin{proof}
The certificate in \lean{Ski/HnfCertData.lean} has 7,244 claims: 7,118 run claims and 126
split claims, in 3,781 strata with ranks at most 2, using 1,860 hints. Every claim checks.
For each of the 2,757 pairs $(p,q)$ of Section~\ref{sec:reduce} there is a claim
on~$X_p\,X_q$. Lemma~\ref{lem:cert} applies.
\end{proof}

\begin{proof}[Proof of Theorem~\ref{thm:main}]
Proposition~\ref{prop:easy} gives acceptance $\Rightarrow$ productivity; Propositions~\ref{prop:reduce} and~\ref{prop:A} give the converse.
\end{proof}

\section*{Mechanization}
All results except the reading of $P_k$ through B\"ohm trees (Section~\ref{sec:prelims}) are
proved in \texttt{lean-ski} version~1.5.0~\citep{leanski}; the statement
of Theorem~\ref{thm:main} is \lean{HnfCert.q5}:
for a closed pure-$S$ term $t$, $(\forall k,\ t\in P_k)\iff c(t)\ne37$. The checker of
Section~\ref{sec:cert} is the Lean function \lean{HnfCert.checkClaim}; its soundness is
\lean{HnfCert.claims\_sound}. Its run on the certificate is checked by
\texttt{decide} in the kernel, without compiled evaluation, split into 642 generated
modules that the build reassembles. The finite facts of Lemma~\ref{lem:table} and
Section~\ref{sec:laws} are checked the same way; the tables are stored as natural-number
literals, so the trusted base includes the kernel's arithmetic on large naturals. Every
theorem is listed in the axiom gate \texttt{Ski/Audit.lean}, which fails the build if a
result depends on an axiom other than \texttt{propext} and \texttt{Quot.sound}.

\bibliographystyle{plainnat}
\bibliography{refs}
\end{document}
